\chapter{Conclusions}
This project focused on designing and implementing a luxury concierge platform for high-net-worth residents and visitors in the Middle East, with a particular emphasis on Dubai. The platform replaces fragmented traditional concierge workflows with a modernized, AI-augmented digital experience built around a tiered membership model. The scope covered mobile application development, backend API design, e-commerce integration, external inventory synchronization, AI-driven itinerary generation, and administrative tooling. By delivering an end-to-end platform that supports browsing, booking, payment, and personalization across both Essential and Elite tiers, this project demonstrates the practical feasibility of combining conversational AI with structured commerce infrastructure to serve a time-poor, high-expectation customer segment.

\section{Key Achievements}

Throughout the development process, the project delivered against the objectives established in Chapter~\ref{chapter:Introduction}. The completed outcomes are:

\begin{enumerate}
  \item A tiered membership platform was implemented with Guest, Essential, and Elite access levels, enforced through Supabase authentication, relational role-permission modeling, and row-level security on sensitive tables.
  \item A unified product inventory was established by combining internally authored experiences with externally synchronized events from the PlatinumList catalog, exposed through a single Medusa-backed API layer consumed uniformly by the mobile client.
  \item A complete booking flow was delivered spanning experience discovery, date and ticket configuration, persistent multi-item cart with reservation hold, Stripe-backed checkout, and post-transaction booking management. The flow supports three distinct pricing models (open-dates, show-time, and rental) through a single adaptive interface.
  \item A conversational AI itinerary builder, ``Ask George'', was built using GPT-4o and the Vercel AI SDK, integrated with a nine-tool function set covering profile access, catalog queries, itinerary assembly, and cart publishing. The builder produces structured, directly bookable itineraries from a single natural-language prompt.
  \item Empirical evaluation of the AI builder was conducted through a five-participant internal study comparing AI-driven and manual itinerary construction, alongside a feature-parity comparison against two AI travel-planning platforms in the broader market.
  \item Administrative tooling for the operations team was delivered through the Medusa admin portal, covering customer management, order and refund management, a sales metrics dashboard, a PlatinumList synchronization extension, and transactional email notifications for key operational events.
  \item A complete CI/CD pipeline was established on Bitbucket Pipelines, with three isolated environments (QA, UAT, Production) provisioned across Vercel, Render, Supabase, and Expo Application Services, enabling tag-triggered promotions and automated database migrations per environment.
\end{enumerate}

\section{Reflections and Lessons Learned}

The delivered platform introduces meaningful improvements over the manual concierge workflows it is designed to augment, both in user-facing capability and in operational infrastructure.

The tiered membership model combined with curated inventory provides a structured alternative to the bespoke-per-request operating model of traditional luxury concierge services, allowing high-expectation users to self-serve for a large set of common requests while preserving a premium experience. The AI itinerary builder, in particular, changes the economics of planning: tasks that traditionally require extended back-and-forth with a human concierge can now be completed in a single conversational session, with the resulting itinerary directly actionable through the same interface. The internal comparative study reinforced this advantage quantitatively, and the feature-parity comparison confirmed that the platform's end-to-end integration between AI planning and booking is not present in competing AI travel-planning products.

Beyond the user-facing features, the project produced a set of supporting capabilities that strengthen the platform's long-term maintainability. The CI/CD pipeline automates build, migration, and deployment across three isolated environments, reducing manual release overhead and improving delivery consistency. The unified inventory abstraction treating internally authored and externally synchronized experiences through a single API surface simplifies downstream consumers and positions the platform to add further external sources without architectural change. The adoption of a monorepo with shared type-safe contracts between the mobile client and the backend API reduces integration risk and accelerates future iteration.

The project also produced organizational learnings that are worth recording. This engagement introduced the team to retrieval-augmented generation and large-language-model integration in a production context for the first time. The learning curve was managed through structured knowledge-sharing sessions and mentorship from a senior engineer on an adjacent team, and the experience reinforced that integrating AI into a commercial product is less about the underlying models themselves and more about the surrounding system design prompt discipline, tool structure, context management, and output contracts. These are lessons that carry forward to future projects in the same space.

\section{Limitation}
\label{section:Limitation}
The Adnexio platform was evaluated through a five-participant internal study (Section~\ref{subsection:Methodology}) and a feature-comparison against AI travel-planning platforms (Section~\ref{subsection:Feature Comparison Against AI Travel-Planning Platforms}). A broader user-satisfaction evaluation with the platform's target demographic was not conducted, and the reasons are worth stating directly.

Target-user access was structurally unavailable during the project period. The platform is designed for high-net-worth individuals in the United Arab Emirates and the broader Middle East region. Conducting user-satisfaction testing with representative members of this demographic was not feasible during the project for two intersecting reasons:

\begin{enumerate}
    \item Regional conditions. The ongoing Middle East conflict during the project period disrupted both the target market's travel and discretionary spending patterns and the client's ability to coordinate user research in the region. Recruiting representative Adnexio users for structured evaluation was not practically achievable under these conditions.
    \item Internal substitution is not a valid proxy. The only available alternative testing with the internal development team or with Thai participants accessible to the project team in Bangkok would not produce findings meaningfully transferable to the actual target demographic. The expectations, spending patterns, and lifestyle contexts of the platform's intended users are sufficiently distinct from any available substitute population that a satisfaction score collected from internal or non-target users would misrepresent the platform's true fitness for purpose.
\end{enumerate}

For these reasons, evaluation was scoped to dimensions that do not depend on representative-user access: speed and accuracy of the AI builder against the platform's own manual planning flow (measurable with any technically capable participant), and feature parity and differentiation against competitor AI travel-planning platforms (measurable by direct observation). These are honest partial evaluations they demonstrate that the system performs as designed and compares favorably on capability against the market without making claims about subjective user satisfaction that the available evidence does not support.

\vspace{1em}

In summary, this project successfully demonstrates the feasibility of combining conversational AI with structured commerce infrastructure to serve a time-poor, high-expectation customer segment. The delivered solution addresses the core pain points of fragmented manual concierge workflows and lays a solid foundation for future expansion to additional vendor integrations, to broader user cohorts, and to deeper personalization as the platform matures. This work reflects the importance of aligning technical decisions with real-world commercial outcomes and highlights the value of systematic planning, adaptive scoping, and honest evaluation in software engineering projects delivered under realistic constraints.

\section{Future Refinements}
\label{section:Future Refinements}

Three improvements would meaningfully strengthen the platform: a smarter AI retrieval layer, a proper user evaluation against the actual target demographic, and a richer set of input modes for the AI builder.

\subsection{Retrieval-Augmented Recommendation via PGVector}

The AI itinerary builder currently retrieves candidate experiences by filtering the Medusa catalog on tags and dates. This works at the present catalog size, but it scales poorly. Every additional product makes the tool-result payload larger, which means more prompt tokens for GPT-4o, slower responses, and higher per-call cost. Tag overlap is also a coarse proxy for what the user actually means in natural language a request for ``a quiet evening with great food'' is poorly served by tag matching alone.

The fix is retrieval-augmented generation. Encode each product's title, description, and metadata as a vector embedding at publish time, and store the embeddings in PostgreSQL using the PGVector extension that the platform already provisions. At query time, embed the user's intent and return only the top-$k$ most semantically similar products instead of a tag-filtered slice. The payload shrinks, response time drops, and recommendations get materially better because semantic intent is captured directly rather than approximated through tags. The same embedding layer also unlocks personalized ordering on the Explore tab later, so the investment compounds beyond the AI builder.

\subsection{Target-Demographic User Evaluation}

The current evaluation was deliberately scoped to what could be measured without access to representative users speed and accuracy of the AI builder, and feature parity against competing platforms. That scoping is methodologically honest, but it leaves the most important question open: does the platform actually serve its intended users well?

The next step is a structured user study with twenty to thirty representative high-net-worth residents and visitors of the United Arab Emirates, conducted once regional conditions and research access allow. The study should measure satisfaction (Net Promoter Score, qualitative interviews), commercial efficacy (conversion from itinerary draft to confirmed booking), and retention (whether users return past the novelty phase). These dimensions cannot be substituted by internal or non-target participants, and they are the dimensions on which any claim about fitness for purpose ultimately rests.

\subsection{Multi-Modal AI Interaction}

The AI builder is text-only today. That under-uses the input modes a luxury concierge experience naturally invites: a spoken request, a photograph of a place the user wants to emulate, a screenshot of an event they have heard about. GPT-4o supports voice and vision natively; the platform's existing AI infrastructure already exposes both message types.

Two upgrades stand out. A voice mode using GPT-4o's realtime audio API would let users dictate a brief and receive spoken responses, matching the white-glove interaction the human concierge offers. Image-grounded queries ``find me something like this'' attached to a photo would add an input mode that none of the competing AI travel-planning platforms currently support, and that maps directly onto how inspiration-driven luxury booking actually happens. Both extend the existing builder rather than replace it, preserving the end-to-end connection from chat to checkout that distinguishes the platform from text-only or generation-only competitors.